%% kapa: exam=1 round=32 year=2021 period=2 form=A subject=accounting subject_ko=회계학 seq=8 number=48
%% source: Q-Net 감정평가사 기출문제 (cst00309) artlSeq=5208397 file=extracted/1008/5208397/A형/2021년도 제32회 감정평가사 1차 2교시 A형.pdf
%% conversion: kordoc 4.24.51 -> tools/gen1.py (auto); review=proofread-sonnet; flags=none
\begin{question}{48}
(주)감평은 20×1년부터 제품판매 ￦5당 포인트 1점을 고객에게 제공하는 고객충성제도를 운영하고 제품판매 대가로 ￦10,000을 수취하였다. 포인트는 20×2년부터 (주)감평의 제품을 구매할 때 사용할 수 있으며 포인트 이행약속은 (주)감평의 중요한 수행의무이다. (주)감평은 포인트 1점당 ￦0.7으로 측정하고, 20×1년 부여된 포인트 중 75\%가 사용될 것으로 예상하여 포인트의 개별 판매가격을 추정하였다. 포인트가 없을 때 20×1년 제품의 개별 판매가격은 ￦9,450이다. 상대적 개별 판매가격에 기초하여 (주)감평이 판매대가 ￦10,000을 수행의무에 배분하는 경우, 20×1년 말 재무상태표에 인식할 포인트 관련 이연수익(부채)은?
\begin{choices}
  \item ￦1,000
  \item ￦1,050
  \item ￦1,450
  \item ￦1,550
  \item ￦2,000
\end{choices}
\end{question}
\answer{1}
